\documentclass[%
  aps, amsmath, amssymb, onecolumn, 11pt
]{revtex4-2}

\usepackage{graphicx}
\graphicspath{{figures/}}
\usepackage{orcidlink}
\usepackage{mathtools}
\usepackage{booktabs}

% -------------- arXiv-Style Margin Stamp ----------------
\usepackage[style=iso]{datetime2}
\usepackage{FiraMono}
\usepackage{tikz}
\usepackage{hyperref}

\AddToHook{shipout/background}{
  \begin{tikzpicture}[remember picture, overlay]
    \node [rotate=-90, color=gray!80, font=\fontfamily{FiraMono-TLF}\selectfont\scriptsize] at ([xshift=-1.2cm]current page.east) {WORKING DRAFT \enspace$\bullet$\enspace SUPPLEMENTAL MATERIAL \enspace$\bullet$\enspace \DTMnow \enspace$\bullet$\enspace \href{https://doi.org/10.5281/zenodo.20279094}{doi.org/10.5281/zenodo.20279094} \enspace$\bullet$\enspace Brian S. Mulloy \enspace$\bullet$\enspace bmulloy@umich.edu};
  \end{tikzpicture}
}
% --------------------------------------------------------

% --- SM-specific counter renumbering ---
\renewcommand{\thesection}{S\arabic{section}}
\renewcommand{\thesubsection}{S\arabic{section}.\arabic{subsection}}
\renewcommand{\thefigure}{S\arabic{figure}}
\renewcommand{\thetable}{S\arabic{table}}
\renewcommand{\theequation}{S\arabic{equation}}

\begin{document}

\title{Supplemental Material for\\``Wigner Resolution from Action Capacity''}

\author{Brian S. Mulloy\orcidlink{0000-0002-1803-3172}}
\email{bmulloy@umich.edu}

\date{\today}

\maketitle

% ======================================================================
% S1 --- A De Gosson--Zurek--Mulloy dictionary
% ======================================================================

\section{A De Gosson--Zurek--Mulloy dictionary}
\label{sec:dictionary}

The Letter fixes one notation. The same phase-space geometry appears in
Zurek's account of sub-Planck structure \cite{zurek2001} and in de
Gosson's symplectic treatment of indeterminacy
\cite{degosson2026indeterminacy, degosson2022polar, degosson2013qblobs,
degosson2005bullsci, degosson2009camel, degosson2024reconstruction},
each under different symbols and different normalizations. This section
maps the Letter's symbols onto theirs. The Letter sets the entries. Zurek
and de Gosson reach well past the Letter, into decoherence, chaos,
Donoho--Stark and Hardy inequalities, Mahler volume, and the John
ellipsoid. Only the part of each work that the Letter's symbols name is
recorded here.

\subsection{Conventions and the convention bridge}
\label{sec:conventions}

The three works measure the same widths with three conventions. The
Letter writes the covariance semi-axes as
\begin{equation}
\Delta x = \sqrt{2}\,\sigma_x, \qquad \Delta p = \sqrt{2}\,\sigma_p,
\label{eq:capacity_semiaxes}
\end{equation}
so that the Heisenberg cell is the state's covariance ellipse and its
area equals the symplectic capacity. Zurek writes the widths as standard
deviations, $L \simeq \sigma_x$ and $P \simeq \sigma_p$, defined from
$L^2 \simeq \langle x^2\rangle - \langle x\rangle^2$ and
$P^2 \simeq \langle p^2\rangle - \langle p\rangle^2$. The two scales
differ by a fixed factor,
\begin{equation}
\Delta x = \sqrt{2}\,L, \qquad \Delta p = \sqrt{2}\,P .
\label{eq:bridge_widths}
\end{equation}
The Letter carries the area factor $\pi$ inside every cell, so the
quantum of action reads $h/2 = \pi\hbar$. Zurek forms bare products,
$A \simeq L \times P$, without the $\pi$. de Gosson measures a cell by
its symplectic capacity $c$, normalized so that a quantum blob has
$c = \pi\hbar = h/2$. Translating the action capacity across the three,
\begin{equation}
A \;=\; \pi\,\Delta x\,\Delta p \;=\; c(\Omega_\Sigma) \;=\; 2\pi\,\sigma_x\sigma_p \;=\; 2\pi\,(L \times P),
\label{eq:bridge_capacity}
\end{equation}
where $\Omega_\Sigma$ is de Gosson's covariance ellipsoid and $L\times P$
is Zurek's action.

One symbol collides across the three works and is worth flagging. In the
Letter $X$ and $P$ are the scaled coordinates $X = x/\Delta x$,
$P = p/\Delta p$, and $X_\theta, P_\theta$ are the rotated quadratures.
In Zurek $P$ is the momentum width and $L$ the position width. In de
Gosson $X \subset \mathbb{R}^n_x$ and $P \subset \mathbb{R}^n_p$ are
convex bodies, the position and momentum data clouds. The dictionary
below keeps each work in its own column so the three uses do not mix.

Zurek fixes these widths only to order of magnitude. He identifies the
support widths $L$ and $P$ with the effective extent of the state,
$L^2 \approx \langle x^2\rangle - \langle x\rangle^2$, forms the bare
action $A \approx L\times P$ without the area factor $\pi$, and writes the
sub-Planck scale as $a \approx \hbar^2/A$. His interference relations
$\delta_x = \hbar/P$, $\delta_p = \hbar/L$, and his tile area, Eqs.~(7),
(8), and~(12), are exact once $L$ and $P$ are set, yet inherit the
order-unity looseness of the support widths. In the Letter's covariance
variables the same relations hold up to the fixed factors of
Eqs.~\eqref{eq:bridge_widths} and~\eqref{eq:bridge_capacity}, $\sqrt{2}$
between the widths and $2\pi$ between the actions. The Letter writes
$\sim$ for this reason and reserves equality for the relations symplectic
polar duality fixes exactly.

\subsection{The dictionary}
\label{sec:table}

Tables~\ref{tab:dict_core} and~\ref{tab:dict_derived} pair each symbol of
the Letter with its counterpart in the two source works. The prose entries
of Sec.~\ref{sec:entries} expand the rows that carry the Letter's central
relations, writing each relation out in all three notations.

\begin{table}[ht]
\centering
\raggedright
\renewcommand{\arraystretch}{1.35}
\begin{tabular}{@{}p{0.30\linewidth} p{0.28\linewidth} p{0.34\linewidth}@{}}
\toprule
\textbf{This Letter (Mulloy)} & \textbf{Zurek \cite{zurek2001}} & \textbf{de Gosson \cite{degosson2026indeterminacy, degosson2022polar, degosson2013qblobs}} \\
\midrule
Position, momentum $x,\ p$ & $x,\ p$ & $x\in\mathbb{R}^n_x,\ p\in\mathbb{R}^n_p$ \\
Wigner function $W(x,p)$ & $W(x,p)$, Eq.~(2) & $W\psi(z)$, Eq.~(33) \\
Capacity semi-axis (position) $\Delta x=\sqrt{2}\,\sigma_x$ & $L$, with $L^2\simeq\langle x^2\rangle-\langle x\rangle^2$ & half-width of $X=\Pi_x\Omega_\Sigma$ \\
Capacity semi-axis (momentum) $\Delta p=\sqrt{2}\,\sigma_p$ & $P$, with $P^2\simeq\langle p^2\rangle-\langle p\rangle^2$ & half-width of $P=\Pi_p\Omega_\Sigma$ \\
Heisenberg cell, action capacity $A=\pi\Delta x\Delta p\ge h/2$ & classical action $A\simeq L\times P$, Eq.~(4) & covariance ellipsoid $\Omega_\Sigma$, $c(\Omega_\Sigma)\ge\pi\hbar$, Thm.~14 \\
Resolution scale (position) $\delta x=\hbar/\Delta p$ & $\delta_x=\hbar/P\simeq\sqrt{2}\,\delta x$, Eq.~(7) & $B_P(\Delta p)^\hbar=B_X(\hbar/\Delta p)$, Eq.~(7) \\
Resolution scale (momentum) $\delta p=\hbar/\Delta x$ & $\delta_p=\hbar/L\simeq\sqrt{2}\,\delta p$, Eq.~(8) & $B_X(\Delta x)^\hbar=B_P(\hbar/\Delta x)$, Eq.~(7) \\
Quantum blob $a_\theta$, action $h/2$ & minimum-uncertainty Gaussian & $QB(S)=S\!\left(B^{2n}(\sqrt{\hbar})\right)$, $c=\pi\hbar$, Eq.~(26) \\
Principal blobs $a_{\theta=0},\ a_{\pi/2}$ & axis interference scales $\delta_x,\delta_p$ & diagonal case $S=M_L$; $\mathrm{John}(X\times X^\hbar)$, Thms.~10, 12 \\
\bottomrule
\end{tabular}
\caption{Dictionary, part 1: core geometry. The Letter's symbols and their
counterparts in Zurek \cite{zurek2001} and de Gosson
\cite{degosson2026indeterminacy, degosson2022polar, degosson2013qblobs}.
Equation and theorem numbers in the Zurek column refer to \cite{zurek2001};
those in the de Gosson column refer to \cite{degosson2026indeterminacy}.
Convention factors relating the columns are collected in
Sec.~\ref{sec:conventions}, Eqs.
\eqref{eq:bridge_widths}--\eqref{eq:bridge_capacity}. Because Zurek
measures with $\sigma$ and the Letter with $\sqrt{2}\,\sigma$, his
resolution scales $\delta_x,\delta_p$ exceed the Letter's
$\delta x,\delta p$ by the fixed factor $\sqrt{2}$. In de Gosson's column
$B_X(R)$ and $B_P(R)$ are balls of radius $R$ in position and momentum
space and $X^\hbar$ is the $\hbar$-polar dual of $X$.}
\label{tab:dict_core}
\end{table}

\begin{table}[ht]
\centering
\raggedright
\renewcommand{\arraystretch}{1.35}
\begin{tabular}{@{}p{0.30\linewidth} p{0.28\linewidth} p{0.34\linewidth}@{}}
\toprule
\textbf{This Letter (Mulloy)} & \textbf{Zurek \cite{zurek2001}} & \textbf{de Gosson \cite{degosson2026indeterminacy, degosson2022polar, degosson2013qblobs}} \\
\midrule
Quorum cell $\tilde a=\pi\delta x\delta p=\pi\hbar^2/(\Delta x\Delta p)$ & sub-Planck action $a\simeq\hbar^2/A$, Eq.~(3); tile $(2\pi\hbar)^2/A$, Eq.~(12) & $\hbar$-polar dual $A^\hbar$; saturated pair $X\times X^\hbar$ \\
Reciprocity $\tilde a\,A=(h/2)^2$, chain $A\supseteq a_\theta\supseteq\tilde a$ & $a\simeq\hbar/N$, $N=A/2\pi\hbar$ & biduality $(X^\hbar)^\hbar=X$, antimonotone, Eqs.~(4),(5) \\
Quantum kernel $K_\theta$ (Wigner function of $a_\theta$) & Husimi smoothing erases $a$ (Methods) & $W\psi_{W,Y}=(\pi\hbar)^{-n}e^{-Gz\cdot z/\hbar}$, Eq.~(41) \\
Convolution $\widetilde W_\theta=W*K_\theta$ & not used & smoothing toward a non-negative distribution \cite{degosson2005bullsci, cartwright1976} \\
Portrait $\widetilde W=\frac{1}{\pi}\!\int_0^\pi \widetilde W_\theta\,d\theta\ge 0$ & not used & not used \\
Quadrature rotation $X_\theta,P_\theta$, angle $\theta$ & N, S, E, W of the compass & symplectic rotation $R\in Sp(n)\cap O(2n)$; covariance $W(\hat S^{-1}\psi)=W\psi\circ S$, Eq.~(34) \\
Scaled coordinates $X=x/\Delta x,\ P=p/\Delta p$; rigid ellipse $\hat a_\theta$ & not used & normalization to $B^{2n}(\sqrt{\hbar})$ via $M_L$ \\
Saturation $A=h/2$ (vacuum), $a_\theta=\tilde a=A$ & smallest scale $a\to\hbar$ & unique blob, $c=\pi\hbar$, saturated $X^\hbar=P$, Thm.~14 \\
\bottomrule
\end{tabular}
\caption{Dictionary, part 2: derived and dynamical quantities. Conventions
and column sources are as in Table~\ref{tab:dict_core}.}
\label{tab:dict_derived}
\end{table}

\subsection{Entry by entry}
\label{sec:entries}

\paragraph{Wigner function.}
The three works share one object. The Letter writes
$W(x,p)=\frac{1}{\pi\hbar}\int \psi^*(x+s)\psi(x-s)\,e^{2ips/\hbar}\,ds$.
Zurek's Eq.~(2) is the same transform with the integration variable
$y=2s$ and the prefactor $1/2\pi\hbar$. de Gosson's Eq.~(33) is the
$n$-dimensional form $W\psi(z)$ on $\mathbb{R}^{2n}$. All three integrate
to one and return the marginals $|\psi(x)|^2$ and $|\widehat\psi(p)|^2$.

\paragraph{Capacity semi-axes and the action capacity.}
The Letter's Heisenberg cell is $A=\pi\,\Delta x\,\Delta p\ge h/2$.
Zurek's classical action is the bare product $A\simeq L\times P$ of the
support widths, Eq.~(4). de Gosson's covariance ellipsoid $\Omega_\Sigma$
carries the same content as a symplectic capacity. His Theorem~14 states
that the quantum condition $\Sigma+\tfrac{i\hbar}{2}J\ge 0$ holds if and
only if $c(\Omega_\Sigma)\ge\pi\hbar$ for every symplectic capacity $c$,
which is if and only if $\Omega_\Sigma$ contains a quantum blob. The
Letter's floor $A\ge h/2$ is this capacity bound, since $h/2=\pi\hbar$
and $c(\Omega_\Sigma)=\pi\,\Delta x\,\Delta p=A$. The conversion to
Zurek's product is Eq.~\eqref{eq:bridge_capacity}, $A=2\pi\,(L\times P)$.

\paragraph{Resolution scales and $\hbar$-polar duality.}
This is the relation the dictionary turns on. The Letter writes
\begin{equation}
\delta x=\frac{\hbar}{\Delta p}, \qquad \delta p=\frac{\hbar}{\Delta x}.
\label{eq:resolution}
\end{equation}
Zurek reaches the same scales for the smallest structures of the Wigner
function, $\delta_x=\hbar/P$ in his Eq.~(7) and $\delta_p=\hbar/L$ in his
Eq.~(8). de Gosson reaches them as $\hbar$-polar duality. For an interval
of half-width $R$ the dual is the interval of half-width $\hbar/R$, his
Lemma~1, Eq.~(7),
\begin{equation}
B_P(\Delta p)^\hbar=B_X\!\left(\frac{\hbar}{\Delta p}\right)=B_X(\delta x),
\qquad
B_X(\Delta x)^\hbar=B_P\!\left(\frac{\hbar}{\Delta x}\right)=B_P(\delta p).
\label{eq:polar_intervals}
\end{equation}
The three statements coincide once the widths are converted by
Eq.~\eqref{eq:bridge_widths}. Because Zurek measures with $\sigma$ and the
Letter measures with $\sqrt{2}\,\sigma$, the Letter's resolution is finer
than Zurek's nominal scale by the same fixed factor,
$\delta x=\delta_x/\sqrt{2}$.

\paragraph{Quantum blob.}
The Letter's $a_\theta$ is a de Gosson quantum blob of action $h/2$,
deforming at fixed action across $\theta$. de Gosson's definition is
$QB(S)=S(B^{2n}(\sqrt{\hbar}))$ for $S\in Sp(n)$, his Eq.~(26), with
volume $\tfrac{1}{n!}(h/2)^n$ and symplectic capacity $c=\pi\hbar$. In one
freedom this is an ellipse of area $\pi\hbar=h/2$, the Letter's $a_\theta$.
Zurek's building block is the round minimum-uncertainty Gaussian, the
coherent-state cell of area of order $\hbar$ from which his cat and compass
states are assembled.

\paragraph{Principal blobs.}
The Letter's two principal members are
$a_{\theta=0}=\pi\,\Delta x\,\delta p=h/2$ and
$a_{\pi/2}=\pi\,\delta x\,\Delta p=h/2$. In de Gosson these are the
diagonal, axis-aligned case. His projection theorem, Theorem~10, gives a
saturated polar pair, $(\Pi_X QB(S))^\hbar=\Pi_P QB(S)$, exactly when
$S=M_L$, and his Theorem~12 identifies the unique blob inscribed in
$X\times X^\hbar$ as the John ellipsoid. Zurek's interference scales
$\delta_x,\delta_p$ are these principal members read off the coordinate
axes.

\paragraph{Quorum cell, polar dual, reciprocity.}
The Letter's quorum cell is
$\tilde a=\pi\,\delta x\,\delta p=\pi\hbar^2/(\Delta x\,\Delta p)$, the
symplectic polar dual of the Heisenberg cell, with the reciprocity
$\tilde a\,A=(h/2)^2$ and the chain $A\supseteq a_\theta\supseteq\tilde a$.
de Gosson's $\hbar$-polar dual $A^\hbar$ is the same construction, with
biduality $(X^\hbar)^\hbar=X$ and antimonotonicity, his Eqs.~(4) and (5).
The reciprocity follows from the interval duals of
Eq.~\eqref{eq:polar_intervals}, since
$\tilde a\,A=\pi^2(\delta x\,\Delta p)(\delta p\,\Delta x)=\pi^2\hbar^2=(\pi\hbar)^2$.
Zurek's sub-Planck action is $a\simeq\hbar^2/A$, his Eq.~(3), with the
compass tile $a=(2\pi\hbar)^2/A$ of his Eq.~(12). In Zurek's widths the
Letter's quorum cell is $\tilde a=(\pi/2)\,\hbar^2/(L\times P)$. He also
records the count $N=A/2\pi\hbar$ of orthogonal states and the scale
$a\simeq\hbar/N$.

\paragraph{Quantum kernel, convolution, portrait.}
The Letter's quantum kernel $K_\theta$ is the Wigner function of the blob
$a_\theta$, a centered Gaussian of action $h/2$. de Gosson's Eq.~(41)
gives the Wigner function of a generalized Gaussian,
$W\psi_{W,Y}(z)=(\pi\hbar)^{-n}e^{-Gz\cdot z/\hbar}$, which is the kernel
written through its covariance. Convolving $W$ with a Gaussian of action
$h/2$ returns a non-negative distribution, the route through Cartwright
\cite{cartwright1976} and de Gosson's Husimi--Wigner cells
\cite{degosson2005bullsci}. The Letter's per-angle convolution is
$\widetilde W_\theta=W*K_\theta$ and the portrait is the family integral
$\widetilde W=\frac{1}{\pi}\int_0^\pi\widetilde W_\theta\,d\theta\ge 0$.
Zurek notes the complementary fact that the Husimi function, the Wigner
function smoothed on a full Planck scale, erases the sub-Planck scale $a$.
The portrait sits between the two, resolved at $\tilde a$ rather than at
$h$.

\paragraph{Quadrature rotation.}
The Letter's rotated quadratures are
$X_\theta=X\cos\theta+P\sin\theta$ and
$P_\theta=-X\sin\theta+P\cos\theta$, the homodyne quadratures at phase
$\theta$ that orient the blob and the kernel. de Gosson carries the same
rotation as a symplectic rotation $R\in Sp(n)\cap O(2n)$ and as the
metaplectic covariance of the Wigner function, $W(\hat S^{-1}\psi)=W\psi\circ S$,
his Eq.~(34). The family integral over $\theta$ is the tomographic sweep
\cite{degosson2022radon, darianomacconeparis2001}. Zurek's compass orients
its lobes north, south, east, and west, fixing the two axis directions of
the interference grid.

% ======================================================================
% Scaffolding for later sections (not rendered).
% The dictionary above is the focus of this revision. The headings below
% are retained as a roadmap for the planned expansion.
% ======================================================================
%
% \section{Action-units table for phase-space methods}
% \section{Effective resolution: derivations}
% \section{Computational details}
%   - pipeline overview, twelve-state list with covariance parameters,
%     integration-grid choices, Zenodo / OSF deposit pointers.

\bibliography{main}

\end{document}